\documentclass[10pt,a4paper]{amsart}
\usepackage[margin=19mm]{geometry}
\usepackage{amsmath,amssymb,mathtools}
\usepackage[scr=rsfs,cal=euler]{mathalfa}
\usepackage{xfrac}
\usepackage[unicode,hidelinks,bookmarks=false]{hyperref}
\newcommand{\ie}{i.e.\ }
\newcommand{\cf}{cf.\ }
\newcommand{\resp}{respectively\ }
\newcommand{\Subsection}{\subsection}
\providecommand{\nobreakdash}{\nobreak\hbox{-}}
\setlength{\emergencystretch}{2em}
\multlinegap0pt
\usepackage{bm}
\usepackage{bbm}
\usepackage{dsfont}
\usepackage{xspace}
\usepackage{cancel}
\usepackage{mathtools}
\usepackage[shortlabels]{enumitem}
\usepackage{amsthm}
\makeatletter
\@ifundefined{theorem}{\newtheorem{theorem}{Theorem}}{}
\@ifundefined{lemma}{\newtheorem{lemma}[theorem]{Lemma}}{}
\@ifundefined{sublemma}{\newtheorem{sublemma}{}}{}
\makeatother
\newcommand{\cS}{\mathcal{S}}
\newcommand{\cl}{{\rm cl}}
\newcommand{\si}{{\rm si}}
\title[The excluded minors of the class of spike minors]{The excluded minors of the class of spike minors}
\author{Sam Bastida, Nick Brettell, Rutger Campbell, George Drummond, Emma Hogan, Charles Semple, Gerry Toft}
\date{Source: arXiv:2511.07570v1 (CC BY 4.0)}
\begin{document}
\maketitle

\noindent\emph{Source and licence.} The excluded minors of the class of spike minors. Version arXiv:2511.07570v1 (CC BY 4.0), distributed under the Creative Commons Attribution 4.0 International licence, as recorded in the arXiv licence metadata for this exact version and shown on the abstract page https://arxiv.org/abs/2511.07570v1. Full rights notice: licenses/arxiv-2511.07570-CC-BY-4.0.txt.\par\smallskip

\noindent\textit{Editorial scope.} Counted unit: the Lemma whose source label is \texttt{lemma: small spike-minors} in the article (source0.tex lines 532--542, source label \texttt{lemma: small spike-minors}), delivered together with the article's own complete original proof of it (source0.tex lines 544--581, one \texttt{proof} environment of 38 lines). No mathematics is written, repaired or re-derived: statement and proof are the article's own text line for line. Required context retained uncounted: lemma: spike degeneracy (lines 455--463); lem:smaller\_spike\_minors (lines 514--522). Every definition, display and label of those lines is retained unchanged. Omitted from this package and not redistributed: 1--454, 464--513, 523--531, 543--543, 582--1161. Nothing inside a retained excerpt is omitted or altered: every retained body is delivered exactly as the article's lines are written, so no omission receipt of any kind accompanies this package. Citations: the retained excerpts carry no citation command at all, so no bibliography entry of the article is transcribed and no reference is redistributed. Reference adaptation: none was needed and none was made; every reference inside the delivered unit resolves inside it, so no unresolved reference or citation is produced. Printed slips kept literal: no printed word, symbol, hypothesis or number of the counted statement or of its proof was altered. Layout and class: the article's own class is \texttt{amsart} with packages including latexsym, amsmath, amsfonts, amssymb, hyperref, mathrsfs, float, subcaption, tikz, euscript, enumerate, cleveref, color, xcolor; the wrapper is the lane's pinned \texttt{amsart} editorial wrapper at 10pt on A4, which restates the theorem-like environments the retained text needs and adds the article's own macro definitions that the retained text needs (\texttt{\textbackslash cS}, \texttt{\textbackslash cl}, \texttt{\textbackslash si}), with extra wrapper packages (bm, bbm, dsfont, xspace, cancel, mathtools, enumitem). The delivered numbering is the unit's own. Assets: no retained span contains a figure environment, a graphics-inclusion command, a PDF-page inclusion, a table or a TikZ picture; no image file is copied into this package, the package declares no asset dependency and no image file is redistributed with it. Licence: version arXiv:2511.07570v1 of this article is distributed under the Creative Commons Attribution 4.0 International licence, as recorded in the arXiv licence metadata for this exact version and shown on the abstract page https://arxiv.org/abs/2511.07570v1. A full rights notice is delivered at licenses/arxiv-2511.07570-CC-BY-4.0.txt. Characters: every retained excerpt is pure ASCII, so the delivered engine, XeLaTeX with its default Latin Modern text font, reports no missing character.
\begin{lemma}
\label{lemma: spike degeneracy}
Let $r\ge 3$, and let $M$ be a rank-$r$ matroid. Then $M$ is a minor of a spike if and only if, for some positive integer $n$, the matroid $M$ is isomorphic to a restriction of either
\begin{enumerate}[{\rm (i)}]
\item $2U_{r, r+1}\oplus U_{0, n}$, or

\item a rank-$r$ tipped spike in which the tip has been replaced by a parallel class of size $n$.
\end{enumerate}
\end{lemma}

\begin{lemma}
\label{lem:smaller_spike_minors}
Let $r\in \{1, 2\}$ and let $M$ be a rank-$r$ matroid. Then $M$ is a minor of a spike if and only if, for some positive integer $n$, the matroid $M$ is isomorphic to a restriction of one of the matroids
\begin{enumerate}[{\rm (i)}]
\item $2U_{2, 3}\oplus U_{0, n}$, $\{n\}$-$U_{2,5}$, $\{2, n\}$-$U_{2,4}$, $\{2, 2, n\}$-$U_{2,3}$,

\item $U_{1, 4}\oplus U_{0, n}$, and $U_{1, n}\oplus U_{0, 1}$.
\end{enumerate}
\end{lemma}

\begin{lemma}
\label{lemma: small spike-minors}
Let $M$ be a matroid with rank at most three. Then $M$ is an excluded minor for $\cS$ if and only if $M$ is isomorphic to one of the matroids
\begin{enumerate}[{\rm (i)}]
\item $U_{1, 5}\oplus U_{0, 2}$,

\item $U_{2, 4}\oplus U_{0, 1}$,\, $U_{2, 6}$,\, $U_{1, 3}\oplus U_{1, 3}$,\, $\{2, 2, 2\}$-$U_{2, 4}$,\, $\{2, 2\}$-$U_{2, 5}$,\, $U_{1, 3}\oplus U_{1, 1}\oplus U_{0, 1}$,

\item $U_{2, 4}\oplus U_{1, 1}$,\, $U_{2, 3}\oplus U_{1, 1}\oplus U_{0, 1}$,\, $\{2\}$-$U_{2, 3}\oplus U_{1, 2}$,\, $P_7^-$, and $P_7^=$.
\end{enumerate}
\end{lemma}

\begin{proof}
Using Lemmas~\ref{lemma: spike degeneracy} and~\ref{lem:smaller_spike_minors}, it is straightforward to verify that each of the matroids in (i), (ii), and (iii) is an excluded minor for $\cS$. It remains to show that there are no further excluded minors for $\cS$ of rank at most three. We check each rank in turn, noting that every rank-$0$ matroid is a spike minor.

\begin{sublemma}
The matroid $U_{1,5} \oplus U_{0,2}$ is the only rank-$1$ excluded minor for $\cS$.
\label{rank1}
\end{sublemma}

Let $M$ be a rank-$1$ excluded minor for $\cS$. Then $M\cong U_{1, n}\oplus U_{0, m}$ for some positive integers $m$ and $n$. If $n\geq 5$ and $m\geq 2$, then $M$ has a minor isomorphic to $U_{1,5} \oplus U_{0,2}$. If either $n < 5$ or $m < 2$, then $M$ is a spike minor by Lemma~\ref{lem:smaller_spike_minors}(ii). Thus (\ref{rank1}) holds.

\begin{sublemma}
The matroids listed in (ii) are the only rank-$2$ excluded minors for $\cS$.
\label{rank2}
\end{sublemma}

Let $M$ be an additional rank-$2$ excluded minor for $\cS$. If $M$ has a loop, then $M$ contains at most three (including trivial) distinct parallel classes; otherwise, $M$ has a minor isomorphic to $U_{2,4} \oplus U_{0,1}$. If every parallel class of $M$ has size at most two, then $M$ is a restriction of $2U_{2, 3}\oplus U_{0, n}$ and so, by Lemma~\ref{lem:smaller_spike_minors}, $M$ is in $\cS$. Therefore, $M$ has a parallel class of size at least three, and so $M$ has a minor isomorphic to $U_{1, 3}\oplus U_{1, 1}\oplus U_{0,1}$, another contradiction. Hence, $M$ has no loops.

Let $\si(M)\cong U_{2, n}$ for some positive integer $n\ge 2$. Since $M$ has no minor isomorphic to $U_{2, 6}$, we have that $n < 6$. Furthermore, as $M$ has no minor isomorphic to $U_{1,3} \oplus U_{1,3}$, $M$ has at most one parallel class containing at least three elements. Hence, if $n \leq 3$, Lemma~\ref{lem:smaller_spike_minors} implies that $M$ is in $\cS$. Therefore $n\in \{4, 5\}$. If $n = 4$, then, as $M$ does not have a minor isomorphic to $\{2,2,2\}$-$U_{2, 4}$, we have that $M$ has at most two non-trivial parallel classes. But then, by Lemma~\ref{lem:smaller_spike_minors}, $M$ is a spike minor, a contradiction. If $n = 5$, then, as $M$ does not have a minor isomorphic to $\{2, 2\}$-$U_{2, 5}$, it follows that $M$ has at most one non-trivial parallel class. But, again by Lemma~\ref{lem:smaller_spike_minors}, $M$ is in $\cS$. This completes the proof of (\ref{rank2}).

\begin{sublemma}
The matroids listed in (iii) are the only rank-$3$ excluded minors for~$\cS$.
\label{rank3}
\end{sublemma}

Let $M$ be an additional rank-$3$ excluded minor for $\cS$. If $M$ has a $U_{2, 4}$-restriction, then $M$ has a minor isomorphic to $U_{2, 4}\oplus U_{1, 1}$, and so $M$ has no such restriction. Suppose $M$ contains neither a triangle nor a parallel class of size at least three. As $r(M) = 3$ and $M$ contains no triangles, we have that $\si(M)\cong U_{3, m}$ for some $m\geq 3$. Since $M$ is not a spike minor, Lemma~\ref{lemma: spike degeneracy} implies that $m\geq 5$. If $M$ has a loop, then $M$ has a minor isomorphic to $U_{3, 5}\oplus U_{0,1}$, and so $M$ has a minor isomorphic to $U_{2, 4}\oplus U_{0, 1}$, a contradiction. Furthermore, if $M$ has a pair of elements in parallel, then contracting one of these elements implies that $M$ has a minor isomorphic to $U_{2, 4}\oplus U_{0, 1}$. So $M$ has no loops or parallel elements. Since $U_{3, 5}$ and $U_{3, 6}$ are both spike minors, $m\geq 7$. But now $M$ has a $U_{2, 6}$ minor, a contradiction. Hence $M$ has either a triangle or a parallel class of size at least three.

Suppose that $M$ has a triangle $T$. If $M$ has a loop, then $M$ has a restriction isomorphic to $U_{2, 3}\oplus U_{1, 1}\oplus U_{0, 1}$. Thus $M$ has no loops. Assume that $M$ has two distinct pairs of parallel elements $\{a_1, a_2\}$ and $\{b_1, b_2\}$ such that $r(\{a_1, a_2, b_1, b_2\}) = 2$. If $\{a_1, a_2, b_1, b_2\}\subseteq \cl(T)$, then $M/a_1$ has a minor isomorphic to $U_{1, 3}\oplus U_{1, 1}\oplus U_{0, 1}$, a contradiction. If $\{a_1, a_2\}\subseteq \cl(T)$ and $\{b_1, b_2\}\not\subseteq \cl(T)$, then $M$ has a minor isomorphic to $\{2\}$-$U_{2, 3}\oplus U_{1, 2}$. Thus we may assume that no element in $\{a_1, a_2, b_1, b_2\}$ is contained in a triangle. But now $M/ a_1$ has a minor isomorphic to $U_{2, 4}\oplus U_{0, 1}$. Hence, $M$ has at most one non-trivial parallel class.

Assume that $M$ has a pair of parallel elements $\{a_1, a_2\}$. Suppose that $\{a_1, a_2\}\not\subseteq \cl(T)$. By Lemma~\ref{lemma: spike degeneracy}(ii), the matroid $U_{2, 3}\oplus U_{1, n}$ is a spike minor for all positive integers $n$. Therefore, there exists an element $e\in E(M) - (T\cup \{a_1, a_2\})$ such that $r(\{a_1, e\})=2$. The matroid $M / a_1$ does not have a minor isomorphic to $U_{2,4} \oplus U_{0,1}$, which implies that there is a triangle $T'$ of $M$ which contains $a_1$, $e$, and an element of $T$. Furthermore, this is true for all elements of $E(M) - (T \cup \{a_1,a_2\})$ which are not in parallel with $a_1$ and $a_2$. Since $M$ has no $U_{2,4}$-restriction, and each of these elements are in a triangle with $a_1$ and an element of $T$, it follows that there are at most three such elements. But this implies that $M$ is a spike minor with the parallel pair $\{a_1,a_2\}$ at the tip and a traversal $T$. This is a contradiction, so $\{a_1,a_2\}$ is in the closure of every triangle of $T$.

Choose a maximal set $\{x_1,x_2,\ldots,x_k\}$ of elements of $E(M) - \{a_1\}$ such that $\{a_1,x_1,x_2,\ldots,x_k\}$ does not contain a non-spanning circuit. Each $x_i$ is either not contained in a triangle, or is contained in a triangle with $a_1$ and an element outside of $\{a_1,x_1,x_2,\ldots,x_k\}$. In particular, if $k \leq 3$, then $M$ is a spike minor with the parallel pair $\{a_1,a_2\}$ at the tip, and each $x_i$ an element from a separate leg. Thus, $k \geq 4$. But now $\{x_1,x_2,x_3,x_4\}$ is a $U_{2,4}$-restriction in $M / a_1$, so $M$ has a minor isomorphic to $U_{2,4} \oplus U_{0,1}$. Hence, $M$ has no parallel elements.

Let $e\in E(M)$. By Lemma~\ref{lemma: spike degeneracy}, the spike minor $E(M \backslash e)$ has at most $7$ elements. If $|E(M)|\le 6$, then it is easily checked that $M$ is a spike minor. Thus $|E(M)|\in \{7, 8\}$. Say $|E(M)| = 8$. Then $e$ is contained in at most three triangles. If $e$ is contained in zero or one triangles, then $M/e$ has a restriction isomorphic to $U_{2, 6}$. If $e$ is contained in exactly two triangles, then $M/e$ is isomorphic to $\{2, 2\}$-$U_{2, 5}$, while if $e$ is contained in exactly three triangles, then $M/e$ is isomorphic to $\{2,2,2\}$-$U_{2, 4}$. Therefore, $|E(M)| = 7$. If $e$ is not contained in a triangle, then $M/e\cong U_{2, 6}$. Hence, every element of $M$ is contained in a triangle. If $M$ has an element common to three triangles, then $M$ is a spike minor. Now, $M$ must have two disjoint triangles, $T$ and $T'$, say. The unique element $f$ of $E(M) - (T \cup T')$ is contained in at least one triangle, and cannot be contained in three triangles. If $f$ is contained in one triangle, then $M \cong P_7^=$, and if $f$ is contained in two triangles, then $M \cong P_7^-$. This is a contradiction, so $M$ has no triangles.

The last case to consider is when $M$ has no triangles, but has a parallel class $\{a_1, a_2, \ldots, a_k\}$, where $k\ge 3$. If $M$ has a second parallel class $\{b_1,b_2,\ldots,b_\ell\}$ with $\ell \geq 2$, then $M / b_1$ has a minor isomorphic to $U_{1, 3}\oplus U_{1, 1}\oplus U_{0, 1}$, a contradiction. Therefore, $\{a_1, a_2, \ldots, a_k\}$ is the only parallel class of $M$ of size at least two. If $|E(M) - \{a_1, a_2, \ldots, a_k\}|\leq 3$, then $M$ is a spike minor. If $|E(M) - \{a_1, a_2, \ldots, a_k\}| > 3$, then, as $M$ has no triangles, $M/a_1$ has a minor isomorphic to $U_{2, 4}\oplus U_{0, 1}$. This last contradiction proves (\ref{rank3}).

Combining (\ref{rank1})--(\ref{rank3}) completes the proof of the lemma.
\end{proof}

\end{document}
